% !TEX root = ../../mixed-sklyanin-cubics.tex
% Current result: 2026-11-23 / v91.
\subsubsection{PBW 濾過による有限収縮と上位 Shamash 成分の構成式}
\label{ifj:subsubsec:filtered-shamash}

本節の有限摂動収縮は基本摂動補題\cite{TWCrainic}の完全複体への適用であり、
高次 homotopy の存在原理は\cite[Definition~3.4.1, Proposition~3.4.2]{CMOEisenbudPeeva}に遡る。
ここでは前節の $H=P_1\to P_2(3)$、$G=P_2\to P_3(3)$ を保持し、
残る一次 homotopy と higher homotopy を\emph{有限行列の逆と積による式}に還元する。
これによって行列全成分を表示した、あるいは比較 Ku 核を得たとは主張しない。
用いる収縮は非可換源の Koszul 複体のもので、
比較対象間の Hom の special retract とは異なる。

\begin{lemma}[一方向 PBW 濾過の有限摂動収縮]
\label{ifj:lem:filtered-koszul-contraction}
許容 $\tau_1,\tau_2$ を固定し、$V=0$, $\Lambda=sU$ とする。
一方向の六段 Koszul 複体 $P_\bullet(s)$ の内部次数 $N>0$ の部分を
$C_\bullet^{(N)}(s)$ と書く。
六つの中心平方と平方なし Clifford 語からなる PBW 基底を用いて
$C_\bullet^{(N)}(s)$ を $C_\bullet^{(N)}(0)\otimes\C[U,s]$ と同一視する。
原点微分を $D_0$、変形微分を $D_s=D_0+\delta_s$ とする。

$C_\bullet^{(N)}(0)$ の $\C$-線型収縮
$h_0:C_i^{(N)}(0)\to C_{i+1}^{(N)}(0)$ を
$D_0h_0+h_0D_0=1$, $h_0^2=0$、かつ二因子の $x$ 次数を保つように選ぶ。
このとき
\begin{equation}\label{ifj:eq:finite-pbw-contraction}
 \boxed{\displaystyle
 h_s=h_0(1+\delta_sh_0)^{-1}
    =h_0\sum_{r=0}^{\lfloor N/2\rfloor}(-\delta_sh_0)^r}
\end{equation}
は $\C[U,s]$-線型な収縮で、
$D_sh_s+h_sD_s=1$, $h_s^2=0$ を厳密に満たす。
特に式\eqref{ifj:eq:finite-pbw-contraction}は形式無限級数ではない。

標準係数 $\tau_1=\tau_2=2$ では、各 $x$ 次数ブロックの
原点微分行列を $D_{i,0}$ とし、通常の有理 PBW 基底における転置を使えば、
\begin{equation}\label{ifj:eq:filtered-koszul-hodge}
 \Delta_i=D_{i+1,0}D_{i+1,0}^{\mathsf T}
              +D_{i,0}^{\mathsf T}D_{i,0},
 \qquad h_{i,0}=D_{i+1,0}^{\mathsf T}\Delta_i^{-1}
\end{equation}
とすることができる。$N>0$ では $\Delta_i$ は有理数上可逆であり、
式\eqref{ifj:eq:finite-pbw-contraction}の係数は有限回の有理数演算だけで指定される。
\end{lemma}

\begin{proof}
分離原点は二つの三生成元 Koszul 複体の tensor 積である。
内部次数 $N>0$ では augmentation の次数零ホモロジーも存在せず、
各 $x$ 次数ごとに完全な有限次元複体となる。
したがって各ブロックで像と補空間を選べば、$h_0^2=0$ を持つ収縮がある。
標準係数の式\eqref{ifj:eq:filtered-koszul-hodge}は有限次元 Hodge 分解である。
複体の完全性と有理実行列性から $\Delta_i$ は正定値、
よって逆行列は有理数成分を持ち、指定恒等式を満たす。

PBW 語の実際の $x$ 次数は中心平方 $x_a^2$ を二つ、$x_a$ を一つとして数える。
一方向の二次関係で $sU$ が出るたび、$x$ 二文字を $y$ 側の中心平方へ置換する。
既存の細分次数 $\deg U_{ab}=(2,-2)$ と、
全 Koszul 基底の細分次数により、$\delta_s$ の $s^r$ 係数は
この\emph{実際の} $x$ 次数を少なくとも $2r$ 下げる。
$h_0$ はその次数を保つため、固定した総内部次数 $N$ では
$(\delta_sh_0)^{\lfloor N/2\rfloor+1}=0$ となる。

$D_s^2=D_0^2=0$ と $D_0h_0+h_0D_0=1$ から
\[
 D_s(1+\delta_sh_0)=(1+\delta_sh_0)D_0
\]
が従う。これと $h_0^2=0$ を使うと
$D_sh_s+h_sD_s=1$, $h_s^2=0$ が直接従う。
この操作は Koszul 複体の有限次元 $\C$-線型操作であって、
$h_s$ 自身の右 $B_s$-線型性は主張しない。
\end{proof}

\begin{proposition}[$P_3\to P_4$ の有限有理行列構成式]
\label{ifj:prop:shamash-fourth-finite-formula}
前二節の $H(s),G(s)$ をそのまま固定する。
$K_3$ の二十基底を $e_\nu$ とし、$N=6$ の式
\eqref{ifj:eq:finite-pbw-contraction}の $C_3^{(6)}\to C_4^{(6)}$ 成分を
$h_{3,s}^{(6)}$ とする。このとき
\begin{equation}\label{ifj:eq:fourth-hpl-formula}
 \boxed{\displaystyle
 J(s)e_\nu
 =h_{3,s}^{(6)}\bigl((w_s\mathbf1_{20}-G(s)d_3(s))e_\nu\bigr),
 \qquad 1\le\nu\le20}
\end{equation}
を各列に適用し、右 $B_s$-線型に延長すると、
\[
 d_4(s)J(s)+G(s)d_3(s)=w_s\mathbf1_{20}
\]
が\emph{全九変数の多項式恒等式}として成立する。
$J(s)$ は $s$ の高々二次であり、$w_y$ 由来の部分は一次以下である。
$\tau_1=\tau_2=2$, $a_1=a_2=1$ で前節の有理数証明書を入力すると、
式\eqref{ifj:eq:fourth-hpl-formula}は、既存の $H,G$ のゲージを保った
$J$ の全係数を有限有理行列逆によって\emph{一意に指定する算法}となる。
ここではその係数を $1120\times315$ の行列として展開・計算実行していない。

原点の作用素 $D_{4,0}:(B_2)^{15}\to(B_3)^{20}$ は
$1120\times315$、階数 $280$ である。
原点の $x$ 次数 $q$ ごとの行列サイズ・階数は
\begin{equation}\label{ifj:eq:fourth-origin-block-ranks}
\begin{array}{c|rrrrrrr}
q&0&1&2&3&4&5&6\\\hline
\dim C_4^{(6)}(q)&0&18&81&117&81&18&0\\
\dim C_3^{(6)}(q)&10&108&270&344&270&108&10\\
\operatorname{rank}D_{4,0}(q)&0&18&72&100&72&18&0
\end{array}
\end{equation}
である。続く一次成分の原点作用素
$D_{5,0}:(B_2)^6\to(B_3)^{15}$ と
$D_{6,0}:B_2\to(B_3)^6$ はそれぞれ
$840\times126$（階数 $120$）、$336\times21$（階数 $21$）となる。
これらの階数は Koszul 完全性からの\emph{理論計算}であり、
指定したピボット小行列の数値検算ではない。
\end{proposition}

\begin{proof}
$d_3G+Hd_2=w_s\mathbf1_{15}$ と $d_2d_3=0$ より
\[
 d_3(w_s\mathbf1_{20}-Gd_3)
 =w_sd_3-(w_s\mathbf1_{15}-Hd_2)d_3=0.
\]
補題\ref{ifj:lem:filtered-koszul-contraction}を各列へ適用すれば
式\eqref{ifj:eq:fourth-hpl-formula}が厳密に解を与える。
$N=6$ では収縮の和は $r\le3$ で停止する。
$x$-曲率の各係数の $x$ 次数は $p+3-p'$、
$y$-曲率では $p-p'$（$p$ は $K_3$、$p'$ は $K_4$ の $x$ 次数）で、
それぞれ最大 $5$、$2$ である。
内部次数二の PBW 係数に $s^mU^m$ が含まれるには
実際の $x$ 次数 $p''\ge0$ に対して $2m+p''$ を必要とする。
従って $m\le2$、$y$-部分では $m\le1$ である。

原点の $B_0$ の二重 Hilbert 級数は
$(1-t_x)^{-3}(1-t_y)^{-3}$ で、
$K_i$ の $x$ 次数 $p$ の基底数は
$\binom3p\binom3{i-p}$ である。
よって
\[
 \dim C_i^{(N)}(q)=
 \sum_p\binom3p\binom3{i-p}
 \binom{q-p+2}{2}\binom{N-i-q+p+2}{2},
\]
ただし負の次数に対応する項を零とする。
$N>0$ の Koszul 完全性と最高段 $P_6$ の単射性から
$\operatorname{rank}D_{4,0}(q)
=\dim C_4^{(6)}(q)-\dim C_5^{(6)}(q)+\dim C_6^{(6)}(q)$。
これを各 $q$ に適用すれば表示の表と階数 $280$ を得る。
同様に $N=7,8$ で $\operatorname{rank}D_{5,0}=120$、
$\operatorname{rank}D_{6,0}=21$ を得る。
\end{proof}

\begin{corollary}[源の $32\mid32$ 曲率模型を指定する有限算法]
\label{ifj:cor:finite-shamash-source-mf}
許容因子を固定した一方向九変数族の Koszul 複体には、
全係数にわたり $D_{\mathrm{MF}}^2=w_s$ を満たす
非可換源の $\Z/2$-次数付き自由 $32\mid32$ 行列分解を、
式\eqref{ifj:eq:finite-pbw-contraction}を用いる有限算法で指定できる。
これは各行列成分を印刷した閉式ではなく、
内部次数 $N\le10$ の有限個の原点行列逆と PBW 積による構成である。
\end{corollary}

\begin{proof}
前命題の $J=s_{1,3}$ に続き、$i=4,5$ で
$R_{1,i}=w_s\mathbf1-s_{1,i-1}d_i$ を列ごとに収縮して
$s_{1,i}:P_i\to P_{i+1}(3)$ とする。
これで $ds_1+s_1d=w_s$ が成立する（$P_6$ の式は $d_6$ の単射性による）。
$s_0=d$ とし、$r=2,3$ に対して
\begin{equation}\label{ifj:eq:higher-shamash-recursion}
 R_{r,i}=-\sum_{a=1}^{r-1}s_{a,i+2(r-a)-1}s_{r-a,i}
          -s_{r,i-1}d_i,
 \qquad s_{r,i}(e)=h_{i+2r-2,s}^{(i+3r)}(R_{r,i}(e))
\end{equation}
と置く。存在する添字だけを用い、$s_{r,-1}=0$ とする。
式の右辺の合成は $P_i\to P_{i+2r-2}$ であり、
既に得た $s_a$ の方程式と中心性から $dR_{r,i}=0$。
各基底に有限収縮を適用し右線型に延長すれば
\[
 ds_r+s_rd=-\sum_{a=1}^{r-1}s_as_{r-a}
\]
を得る。$r=2$ の対象は $0\le i\le3$、
$r=3$ では $0\le i\le1$ に限られる。
残る最上段の恒等式は $d_6$ の単射性から従う。
$r=4$ の障害 $s_1s_3+s_2^2+s_3s_1$ は
$P_0\to P_6$ の成分だけを持ち、$d_6$ との合成が零なので零である。
$r\ge5$ の障害は長さ六から自動的に零になる。
したがって $D_{\mathrm{MF}}=d+s_1+s_2+s_3$ の偶奇折畳みは
$D_{\mathrm{MF}}^2=w_s\mathbf1$ を満たし、
自由rankは $1+15+15+1=32$ と $6+20+6=32$ になる。
これは内部次数を忘れた偶奇折畳みであり、指定 rotation・三乗周期を保つ
graded/window 模型への昇格はここで主張しない。
\end{proof}

\begin{remark}[証明範囲と次の計算]
式\eqref{ifj:eq:higher-shamash-recursion}は\emph{有限線型代数による厳密な指定}であるが、
今回 $J,s_2,s_3$ の全 PBW 行列を有理数で評価したわけではない。
原点の指定 $S_3\times S_3$ 外積模型との厳密なゲージ同定、
高次 homotopy の群同変化、完備な係数証明書も未実行である。
特に、この源の $32\mid32$ 行列分解は実 Ku 比較核、
その Hom の special retract $(i,p,h)$、unit/counit、
rotation と指定三乗周期の coherence を与えない。
固定 Hesse 切片の四十六原式係数・十二比較係数の数値も未計算である。
次は $N=6$ の $280$ ピボットを実際に指定し、
式\eqref{ifj:eq:fourth-hpl-formula}の九変数係数と
非可換 $20\times20$ 恒等式を独立の厳密検算器で照合する。
\end{remark}
